
% SGA 3，报告 VII_B；英文重建本第一轮的中文译本。
% 范围： 节 4.1.2 全文 依据文献局部页 75;
% 止于下节之前： 节 4.2.
% 合并阅读本页： 588. 重版印刷页： 577.
% 依据： Polo--Gille Exp7B-13oct24.pdf,
% SHA-256 D303690E0195654BBE7C3DAA5D89EA23E24CD709A3575217234AC19AE8836694.
% 参照了 Jacob Reinhold 的 CC BY 4.0 Markdown，作为对照来源。
% 本组件无图；所有数学内容均为原生 LaTeX。
\subsubsection*{4.1.2}
\phantomsection
\label{sga3:VIIB:4.1.2}
设 \(G\) 是形式 \(k\)-群。由于函子
\[
  X\longmapsto X^{(p)}
  =
  X\widehat\otimes_\pi k
\]
与有限积交换，它将形式 \(k\)-群映为形式 \(k\)-群。此外，\(\operatorname{Fr}(X/k)\) 关于 \(X\) 是函子性的，所以
\[
  \operatorname{Fr}:G\longrightarrow G^{(p)}
\]
是形式 \(k\)-群同态。
递归定义
\[
  G^{(p^{\,n+1})}
  =
  \bigl(G^{(p^n)}\bigr)^{(p)}.
\]
则复合
\[
\begin{aligned}
  \operatorname{Fr}^{n}
  =
  \operatorname{Fr}^{n}(G/k):
  G
  &\xrightarrow{\ \operatorname{Fr}\ }
  G^{(p)}
  \xrightarrow{\ \operatorname{Fr}\ }
  G^{(p^2)}
  \xrightarrow{\ \operatorname{Fr}\ }
  \cdots
  \xrightarrow{\ \operatorname{Fr}\ }
  G^{(p^n)}.
\end{aligned}
\]
也同样是形式群同态。
令 \(A\) 为 \(G\) 的仿射代数，\(I_A\) 为其增广理想。
\medskip
\noindent\textbf{定义。}
\begin{enumerate}
\item[\textup{(a)}] 将 \(\operatorname{Fr}^{n}\) 的核记为
\[
  \operatorname{Fr}_nG.
\]
由定义直接可得，\(\operatorname{Fr}_nG\) 是无穷小的，其仿射代数为
\[
  A/I_A^{\{p^n\}},
\]
其中 \(I_A^{\{p^n\}}\) 是 \(A\) 中由 \(I_A\) 的元素的 \(p^n\) 次幂生成的闭理想。
\item[\textup{(b)}] 若
\[
  I_A^{\{p^n\}}=0;
\]
则称群 \(G\) 的\emph{高度至多为 \(n\)}；等价地，
\[
  \operatorname{Fr}_nG=G.
\]
\end{enumerate}
